\documentclass{article}
\usepackage[utf8]{inputenc}
\usepackage[T1]{fontenc}
\usepackage[margin=2cm]{geometry}
\usepackage{hyperref}
\sloppy
\begin{document}
\section*{Historical source}
Encyclopaedia Britannica scan, source row 1324. Parse 3.5 extracted the text. Super 3.5 VL corrected it from an enlarged crop of each block. It can still contain recognition errors.
Page headers, footers, and page numbers are removed.

\subsection*{Untitled continuation (E0108)}
\textit{Continues an entry from an earlier page.}

Hedges. price of lime is high they may be built without any cement, and answer the purpose very well if the work is properly executed.”

In making a new fence of this kind, the surface of the ground should be pared off the breadth of the ditch, and likewise for two feet more, in order to prevent as much as possible the thorns from being injured by the growth of grass and weeds. The ditch should be five feet broad, two and a half in depth, and one foot broad at the bottom. Leave one foot for an edging or fencing, then dig the earth one foot of a spade for about one foot, and put about three inches of good earth below the thorn, which should be laid nearly horizontal, but the point rather inclining upwards, in order to let the rain drip to the roots; then add a foot of good earth above it; leave three or four inches of a fencing before another thorn is planted; it must not be directly over the lower one, but about nine inches or a foot to one side of it; then throw a foot of good earth on the thorn, and trample it well down, and level the top of the bank for about three feet and a half for the base of the wall to rest on. This base should be about nine or ten inches, but must not exceed one foot from the thorn. The wall ought to be about two feet thick at the bottom and one foot at the top; the cope to be a single stone laid flat; then covered with two sods of turf, the grass of the undermost to be next the wall, and the other sod must have the grass side uppermost; the sods should be of some thickness, in order to retain moisture; so that they may adhere together, and not be easily displaced by the wind; the height of the wall to be two feet and a half, exclusive of the sods; which together should be from four to six inches, by which means the wall would be near to three feet altogether. The expense of the fences cannot easily be counted, on account of the difference of the prices of labour in different parts. Mr Erskine had them done with lime, every thing included, from 10d. to 13d. per ell (which is equal to 37 inches 2 parts), according to the ease or difficulty of working the quarry, and the distance of it from the place where the fence is erected. The lime costs about 6d. per boll of about 4082667 bolls; and from 15 to 16 bolls of lime are used to the rood of 36 square ells Scots measure; and there are upwards of 43 Scotch ells, or 44 English yards. When the common round or flint stones are made use of, as they require more lime, it is necessary to use 30 or 35 bolls of lime to the rood. The thorns are sold from five to ten shillings per thousand, according to their age, reckoning six score to the hundred. Making the ditch, laying the thorns, and preparing the top of the wall, generally cost from 7d. to 8d. every six ells. About 50 carts of stones, each cart carrying from 7 to 9 cwt. will build a rood; the carriage at 2d. per cart for half a mile distance.

Warmth is undoubtedly as beneficial to hedges as trees; and the walls give an effectual shelter, which in exposed situations is absolutely necessary for rearing young hedges; and they likewise prefer a proper degree of moisture about the roots. If the hedges have been planted for six or seven years before the wall is built, cut them over to two or three inches above the ground with a sharp tool, either in October or November, or early in the spring; and erect the

wall as quickly in that season as possible (the spring in this country can scarcely be said to begin till the end of March). It is almost impossible to imagine the rapidity with which hedges grow in favourable fituations. Mr Erftine had one cut over in the spring, and by the end of the year it was almost as high as the wall. In three years he supposed, that not even the Highland sheep, who easily overleap a wall of four feet and half in height, would have been able to break through it.

Notwithstanding the reasons that have been given already against the planting of timber trees in hedges, we find the practice recommended by some authors as one of the best situations for raising hip-timber. The reasons are, that the roots have free range in the adjoining inclosures, and the top is exposed to the exercise of the winds; by which means the trees are at once enabled to throw out strong arms, and have a large spreading head at the same time; so that we thus at once obtain quickness of growth, with strength and crookedness of timber. Well trained timber trees it is alleged are not prejudicial to hedges, though pollards and low spreading trees are destructive to the hedge-wood which grows under them; neither are high trees prejudicial to corn-fields like high hedges and pollards, which prevent a proper circulation of air; and in Norfolk, where the cultivation of grain is carried on in great perfection, such lands are said to be wood-bound. But when an hedge is trimmed down to four or five feet high, with oaks interspersed, a circulation of air is rather promoted than retarded by it; and a trimmed hedge will thrive quite well under tall flemmed trees, particularly oaks. For arable inclosures, therefore, hedges are recommended of four or five feet high, with oak-timbers from 15 to 25 feet stem. Higher hedges are more eligible for grass lands: the grasses affect warmth, by which their growth is promoted, and consequently their quantity is increased, though perhaps their quality may suffer some injury. A tall fence likewise affords shelter to cattle, provided it be thick and close at the bottom; but otherwise, by admitting the air in currents, it does rather harm than good. The shade of trees is equally friendly to cattle in summer: for which reason it is recommended in grass inclosures to allow the hedge to make its natural hoots, and at the same time to have oak-trees planted in it at proper intervals. Upon bleak hills, and in exposed situations, it will be proper to have two or even three rows of hedgewood, about four feet distant from each other; the middle row being permitted to reach, and always to remain at, its natural height; whilst the side rows are cut down alternately to give perpetual security to the bottom, and afford a constant supply of materials for dead hedges and other purposes of underwood.

Much has been said of the excellence of the holly as a material for hedges; and indeed the beauty of the plant, with its extreme closeness, and continuing green throughout the winter, evidently give it the preference to all others; and could it be raised with equal ease, there is no doubt that it would come into universal practice. Besides the above properties, the holly will thrive almost upon any soil; but thin-soiled stony heights seem to be its natural situation; and it may probably enough be said, that holly will grow


\section*{Source record}
\begin{description}
\item[Scan ID] \texttt{britannica:309f765b1ecc83e0a7f7762e44b4a927a7555443:train:1324}
\item[Image file] \texttt{page-1324-2cfbaa25a150.jpg}
\item[Image SHA-256] \texttt{2cfbaa25a150c2ea23ac854aafdb05fa22f350334e1404f609866a252d061217}
\item[Dataset] NationalLibraryOfScotland/encyclopaedia\_britannica\_illustrated at revision \texttt{309f765b1ecc83e0a7f7762e44b4a927a7555443} (CC0-1.0)
\item[Extraction] model \texttt{nvidia/nemotron-3.5-parse}, raw reply \texttt{extract/raw/page-1324-2cfbaa25a150.json}
\item[Raw reply SHA-256] \texttt{1b39606b6b088a1f85101162677525a11aeb9eec27871f93da22d141e231ce80}
\item[Correction] model \texttt{nvidia/Nemotron-3.5-Super-VL}, block crops, raw replies \texttt{correct/raw/crop\_nothink/page-1324-2cfbaa25a150-b*.json}; super\_vl: 6
\item[Running header] none
\end{description}
\end{document}
